\section{Chapter 2}

\begin{quotex}
\begin{verse}
For everything under this heaven, the beautiful conceived\\
The ugly is born (as correlative)\\
The good established\\
The not-good takes form.\\
Likewise: being and non-being mutually condition each other\\
The possible and the impossible are complementary differentiations\\
Large and small characterize each other\\
The high flips over to the low\\
Clear sounds and noise complement each other\\
Before and after or: ahead and behind follow each other in a circle\\
The True Man is like that\\
He endures in non-acting\\
He teaches without speaking\\
He directs without touching without commanding \\
He forms he makes happen, he leads to development without taking control\\
He accomplishes without doing without calling attention to the one who acts \\
Essentially: not residing in the domain of the correlative where the game of opposites unfolds \\
He always participates in the original force.
\end{verse}
\end{quotex}

\paragraph{Commentary}


In the first lines, referring on one hand to the law of correlation of all human notions; then, on the other hand, to its ontological counterpart, to that which can be called the ``dialectic of the real''. Taoistically: the solidarity and action alternating as opposites, Yin and Yang, gives rise to modifications of the phenomenal world that now complete each other, then pass one into the other.

This dynamism, previously in the \emph{I Ching} given by the circle of the ``Figures of Heaven'', that is, the trigrams and hexagrams, does not touch the calm and undifferentiated essence of the Principle. Chuang tze: ``The immobile center of a circle on the circumference of which crowns all contingencies, distinctions and individuals''. From this, it describes the non-acting behavior of the True Man, of the \emph{sheng jen}, insofar he reflects in himself the Principle in its aspect without name.

On the gnosiological plane, the relativity and the irrelevance of human distinctions is deduced from the principle of correlation of opposites; therefore, of all the current values which correspond to nothing in reality: ``only great spirits are capable of understanding.'' (Chuang Tze)

\paragraph{Chinese text and literal translation}


Chapter 2 (\begin{CJK*}{UTF8}{bsmi}第二章\end{CJK*})
\columnratio{.35}
\begin{paracol}{2}
\begin{CJK*}{UTF8}{bsmi}
天下皆知美之爲美，

斯惡已；

皆知善之爲善，

斯不善已。

故

有無相生，

難易相成，

長短相形，

高下相傾，

音聲相和，

前後相隨。

是以

聖人處無爲之事，

行不言之敎，

萬物作焉而不辭，

生而不有，

爲而不恃，

功成而弗居。

夫唯弗居，是以不去。
\end{CJK*}
\switchcolumn
\vspace*{-2ex}
\begin{verse}
If everybody knows what beauty is,\\
then beauty is not beauty anymore;\\
If everybody knows what goodness is,\\
then goodness is not goodness anymore.\\
Therefore,\\
Being and nothing give birth to one another,\\
Hard and easy are mutually formed,\\
Long and short shape each other,\\
High and low complement each other,\\
Music and voice are harmonized with each other,\\
Front and back follow one another.\\
Hence,\\
The sage focuses on non-action in his works,\\
Practices not-saying in his speech,\\
The myriad things arise but are disregarded\\
The sage produces but does not own\\
Acts but does not claim\\
Accomplishes work but does not focus on it\\
Does not focus on it, and thus it does not go.
\end{verse}
\end{paracol}
\flrightit{Posted on 2020-03-21 by Lao Tzu}

